% "World of Adventure" campaign guide - Combat
% Written by Christopher Thomas.

\chapter{Combat}
\label{sect-combat}

%
%
\section{The Big Changes}

%
\subsection{Armour Rating}

Armour has an ``armour rating'', rather than directly granting an AC bonus.
Shields likewise have an armour rating rather than an AC bonus.

A character's AC is half their armour rating rounded down, and a
character's DR is half their armour rating rounded up.

Per the spell descriptions, Mage Armour and Shield also give armour bonuses
instead of AC bonuses. ``Dodge'', ``deflection'', and ``luck'' bonuses to
AC tend to remain AC bonuses.

Some types of damage (such as sonic damage and damage from adamantine
weapons) can bypass armour-based DR.

%
\subsection{Vigor and Wounds}

Instead of ``hit point'', characters have ``vigor points'' and ``wound
points''.

Vigor points are gained at level-up in the same manner as hit points. These
represent superficial cuts, bruises, and other injuries that are not
life-threatening.

PCs get maximum vigor points at first level, and after that can either roll
(re-rolling any result of ``1'') or can take half-max-plus-one (the
average roll rounded up).

Wound points do not change with level. Damage to wounds represents
life-threatening injuries. Wound points are equal to a creature's
constitution score multiplied by a size multiplier (modified from the SRD's
values):

\begin{tabular}{lc} \hline
\textbf{Size} & \textbf{Multiplier} \\ \hline
Fine & $\times \frac{1}{8}$ \\
Diminutive & $\times \frac{1}{4}$ \\
Tiny & $\times \frac{1}{2}$ \\
\hline
Small & $\times 1$ \\
Medium & $\times 1$ \\
Large & $\times 1\frac{1}{2}$ \\
\hline
Huge & $\times 2$ \\
Gargantuan & $\times 4$ \\
Colossal & $\times 8$ \\
\hline
\end{tabular}

Some types of damage (such as fall damage or natural 20 critical hits)
goes straight to wounds. Most forms of damage deplete vitality first.

Wound points are not affected by \textit{temporary} changes in constitution.
This includes changes caused by spells that change size (Enlarge, Reduce) or
that directly affect constitution (Bear's Endurance), changes due to
temporary ability score damage from poisons, and changes due to creature
attacks that deal temporary ability score damage.

At zero wound points, a character is disabled. At negative wound points, a
character is dying until stabilized. When wound points reach a negative
value the same size as its normal positive maximum, the character dies.

%
\subsection{Critical Hits}

There are two types of critical hit: Hits that are within a weapon's threat
range, and natural 20 critical hits. Neither of these requires
``confirmation''.
\begin{itemize}
\item A natural 20 bypasses DR and goes straight to wounds. It is \textit{not}
multiplied.
\item A threat that was not a natural 20 goes to vitality and \textit{is}
affected by DR, but does more damage (usually twice the damage). Dice are
rolled twice and bonuses are applied twice.

This combines additively, not multiplicatively, with other weapon effects
that multiply damage (per the ``Lance'' example below).
\end{itemize}
Critical hits work on anything that could reasonably have vulnerable
points, including corporeal undead.

A natural 2 is a critical failure. The attack misses and something bad
happens. The exact outcome is at the DM's discretion (typically dropping
or throwing the weapon, tripping and falling prone, or granting nearby
opponents an attack of opportunity).

Using 2d10 instead of 1d20 alters critical threat ranges:
\begin{itemize}
\item Critical threat ranges in the SRD are widened by 1; for example, a
threat range of ``19-20'' becomes ``18-20''.
\item Anything that would double a critical threat range instead widens it
by 1.
\end{itemize}

%
\subsection{Nonlethal Damage}

Nonlethal variants of some weapons and spells exist, per Sections
\ref{sect-equip} and \ref{sect-magic}. These typically have a damage die one
size smaller than the lethal variants.

Lethal and nonlethal damage both affect vitality the same way; the difference
between lethal and nonlethal damage is how they are applied to wound points.

Once vitality points are exhausted, nonlethal damage to wounds is tracked as
its own tally rather than applied to wound points. When nonlethal damage to
wounds is equal to or greater than the \textit{current} number of wound
points (maximum minus lethal damage), a character is unconscious. When
nonlethal damage to wounds equals the \textit{maximum} number of wound
points, any further nonlethal damage to wounds is lethal damage.

%
%
\section{Other House Rules}

%
\subsection{Adamantine}

In addition to being nearly invulnerable, adamantine has a supernatural
ability to cut through almost any material.

In mechanical terms, a slashing or piercing weapon made from adamantine
treats all armour as having an armour rating of zero (no contribution to AC
or to DR). The sole exception is force effect armour (from spells such as
Mage Armour and Shield).

Artificial armour that is in the way of an attack receives damage, taking
half damage from slashing weapons and one point per die from piercing
weapons. Armour that is the direct target of an attack (such as from a
``sunder'' attack) takes full damage from slashing and half damage from
piercing adamantine weapons.

For attacking objects (including armour and weapons via a ``sunder'' attack),
adamantine slashing and piercing weapons bypass hardness.

Adamantine bludgeoning weapons could be built but do not confer a
mechanical advantage to attacks.

%
\subsection{Fall Damage}

Fall damage goes straight to wounds. Damage reduction from armour
\textit{does} apply. ``Suddenly not flying'' is very dangerous even for
high-level characters and for powerful creatures like dragons and wyverns.

Special circumstances (such as landing in a hay cart, or in water, or on to
deep snow) are adjudicated by the DM. This may reduce the effective distance
fallen or provide DR or both at the DM's discretion.

Fall damage is always rolled with d6 dice. The maximum number of damage
dice, and the distance increment per die (rounded down), change with
creature size. The practical upshot is that a squirrel and a dragon both
handle a 20-foot fall better than a human would.

\begin{tabular}{lcc} \hline
\textbf{Size} & \textbf{Max Dice} & \textbf{Distance Per Die} \\ \hline
Fine & 1 & 30' \\
Diminutive & 2 & 20' \\
Tiny & 5 & 10' \\
\hline
Small & 10 & 10' \\
Medium & 20 & 10' \\
Large & 30 & 10' \\
\hline
Huge & 40 & 20' \\
Gargantuan & 60 & 40' \\
Colossal & 80 & 80' \\
\hline
\end{tabular}

NOTE - For 10 dice or more, the damage dealt will be very close to the
average, and tallying dice will be quite cumbersome. Parties that
encounter long falls can optionally use one of the following alternative
approaches:
\begin{itemize}
\item Roll one die, and multiply by the die count. This gives highly
variable damage, which might be the group's preference.
\item Take the flat average (3.5 points per die). This has the opposite
problem (no chance of getting lucky and surviving a long fall).
\item Choose some smaller number of dice (such as 2 or 3), roll and multiply
that, and then roll any remaining dice individually. This gives a bell
curve but still allows variation. The downside is that it's complicated.
\end{itemize}

%
\subsection{Flying Mounts}

Flying mounts are usually at least one size larger than their passengers.

A flying creature can carry up to 10\% of its weight in passengers or
cargo without penalty. Carrying more than that requires constitution checks
at regular intervals. The first check is DC~10; each subsequent check before
the next long rest increases the DC by 1:

\begin{tabular}{lll} \hline
\textbf{Load} & \textbf{Interval} & \textbf{Recovery Time} \\ \hline
up to 10\% & no check needed & n/a \\
up to 20\% & every hour & 1 hour \\
up to 50\% & every 10 minutes & 10 minutes \\
up to 100\% & every round & 1 minute \\
$>$ 100\% & cannot fly & n/a \\ \hline
\end{tabular}

After a failed check, the flying creature gains the ``fatigued'' condition
and must descend until it lands. It is no longer able to maintain level
flight or climb while carrying its load.

The creature must rest for the indicated recovery time before attempting to
fly again (even without load). If made airborne by any means before that
time, it must descend and begin recovering again.

After recovering, the creature can try picking up the load again and flying
with it again, but the DC does not decrease until after the next long rest.
Since the creature is still fatigued due to failing the first check, a
subsequent failed check results in the creature gaining the ``exhausted''
condition.

An ``exhausted'' creature cannot take off or maintain level flight, and must
descend and land per above if made airborne by any means.

%
\subsection{Grab Attack}

Creatures that have the ``grab'' special ability with a bite attack can
automatically apply damage on subsequent rounds if they succeed at grappling
their target (mechanically the same as the ``constrict'' ability).

If a creature has part of your body in its mouth, chewing on you is free.

At the DM's discretion this may apply to other natural attacks, adjudicated
on a case-by-case basis.

%
\subsection{Grappling}

I am modifying how grappling between creatures of different sizes works:
\begin{itemize}
\item \textit{In addition} to the existing size bonuses to CMB and CMD, the
larger creature gets a +5 bonus to CMB and CMD if there is a 1-step size
difference, and a +10 bonus if the size difference is 2 steps or greater.
\item A creature 3 or more size steps larger than another never gains the
grappled condition from the smaller creature.
\end{itemize}

I am also adjusting the rules for grappling someone without being grappled
in return:
\begin{itemize}
\item If a creature wants to grapple without being grappled in return, the
attacker takes a -10 penalty to CMB for making the grapple, and the defender
gets a +10 bonus to CMB when attempting to escape the grapple.
\item If the defender is pinned, the penalty and bonus are 5 rather than 10.
\item This supersedes the SRD's rules (which typically used a penalty of 20).
\item If you creatures are grappling each other, the creature in control of
the grapple can convert it to ``grapple without being grappled in return''
by making a grappling check with the above penalty.
\end{itemize}

The practical effects are as follows:
\begin{itemize}
\item An adult can pick up and carry a child without being counted as
grappled by the child.
\item A flying creature can snatch and carry off a target without an absurdly
high penalty (just a net -5 for a wyvern carrying off a human, for example).
\item It is easier grapple and pin someone and then convert it to restraining
them without being grappled than it is to go straight for the restraint.
\item Mr. Welch can no longer ``abuse the grappling rules with mobs of
trained squirrels'', since squirrels are too small to give a human the
``grappled'' condition. Doing it with chihuahas would be difficult too, since
they'd all be at a -10 penalty.
\end{itemize}

%
\subsection{Lances}

Per the SRD, a weapon such as a lance that deals double damage under certain
conditions will roll twice as many dice \textit{and} add any relevant bonuses
twice.

This combines additively, not multiplicatively. For example, a critical hit
with a lance during a mounted charge does 3 times the base damage, not 4
times.

%
\subsection{Power Attack}

Creatures with secondary attacks only get half the damage bonus from Power
Attack on those secondary attacks. Primary attacks get the full bonus.

%
\subsection{Sonic Damage}

Sonic damage treats all armour as having an armour rating of zero (no
contribution to AC or DR).

Spells that deal sonic damage are typically one level higher than similar
spells that deal other types of damage.

%
\subsection{Start of Combat}

Unless you're surprised, you aren't flat-footed at the start of combat.
Under normal conditions, both sides have enough awareness of the situation
to get full AC and to take a normal round's worth of actions.

%
% This is the end of the file.
